\chapter{Conclusion and Future Work}
\label{ch:conclusion}

\section{Summary of contributions}

This thesis designed, implemented, and evaluated \padonap{}, a
four-stage AI-augmented NFV orchestration pipeline for proactive DDoS
mitigation in data-center networks. The four technical contributions,
each backed by verifiable artefacts, are:

\begin{enumerate}
  \item \textbf{Architecture} (\texttt{pipeline/}, \texttt{graphify-out/}).
        A four-stage pipeline (S1 telemetry → S2 features → S3 AI →
        S4 orchestration) with a canonical \texttt{AIOutputPayload} JSON
        interface on DMaaP, keeping the AI and orchestration layers
        independently replaceable. The architecture incorporates a
        five-tier graduated policy engine with hysteresis and a
        frequency guard to prevent VNF churn.

  \item \textbf{AI components} (\texttt{pipeline/s3\_ai/}, \texttt{models\_v2/}).
        An XGBoost 7-class detector (Acc = 0.9825, Macro-F1 = 0.9642,
        AUC = 0.9958, P99 = 2.2~ms) combined with a Transformer+LSTM
        4-horizon forecaster (AUC $0.988 \to 0.971$ over
        $t{+}30 \to t{+}120$~s, P99 = 7.68~ms). Both trained with
        temporal split, validated with three CV strategies, and
        augmented with FGSM adversarial examples.

  \item \textbf{LP multi-tenant SLA allocator}
        (\texttt{pipeline/s4\_orchestration/sla\_allocator.py}).
        A HiGHS-backed linear program that guarantees per-tenant
        bandwidth floors under VNF overhead up to 700~Mbps, with a
        proportional fallback. URLLC floor (150~Mbps) preserved in
        all 8 evaluation scenarios.

  \item \textbf{Quantitative evaluation} (\texttt{evaluation/}).
        8/8 AI scenarios pass; threshold baseline passes only 6/8
        (fails S4/S5 OOD via over-escalation). Lead-time advantage:
        +65~s on gradual SYN-ramp (S3), +135~s within-orchestrator
        on novelty scenario (S8). Proactive $T_2$ latency 505~ms
        vs reactive $T_3$ latency 6006~ms ($10.9\times$ speed-up).
\end{enumerate}

%%─────────────────────────────────────────────────────────────────────────────
\section{Hypothesis attainment}
\label{sec:kpi}

Table~\ref{tab:hypotheses} maps each hypothesis from Chapter~\ref{ch:intro}
to the experimental evidence in Chapter~\ref{ch:eval}.

\begin{table}[H]
\centering
\small
\caption{Hypothesis attainment summary.}
\label{tab:hypotheses}
\begin{tabular}{@{}lllll@{}}
\toprule
ID & Hypothesis & Target & Attained & Source \\
\midrule
H1 & Lead-time on gradual attack & $\geq 30$~s & \textbf{+65~s (S3)} &
     Table~\ref{tab:lead-time} \\
H2 & Detection Acc / AUC & $\geq 0.98$ / $\geq 0.99$ & \textbf{0.9825 / 0.9958} &
     Table~\ref{tab:ai-metrics} \\
H3 & URLLC floor preserved & $\geq 150$~Mbps, all tiers &
     \textbf{150~Mbps, 8/8 scenarios} & \S\ref{sec:sla-fairness} \\
\bottomrule
\end{tabular}
\end{table}

All three hypotheses are supported by experimental evidence.

%%─────────────────────────────────────────────────────────────────────────────
\section{Limitations}

We summarise the five threats to validity declared in
\S\ref{sec:validity}:
(i) VNF boot times simulated from ETSI PoC calibrations, not live
Docker measurements;
(ii) ONAP SO runs as a stub implementing the v7 REST contract, not
a live ONAP OOM deployment;
(iii) three of seven attack classes absent from CICDDoS2019 training,
evaluated only as OOD;
(iv) single-node evaluation, no distributed multi-orchestrator test;
(v) single-dataset, no cross-dataset generalisation study.
These are disclosed limitations, not concealed.

%%─────────────────────────────────────────────────────────────────────────────
\section{Future work}
\label{sec:future}

Five directions follow directly from the limitations above.

\begin{enumerate}
  \item \textbf{Real ONAP OOM integration.} Deploy SO + CLAMP + Policy
        on a Kubernetes cluster; set \texttt{PAD\_ONAP\_STUB=false}
        and replay scenario S8 end-to-end. Measure actual VNF boot time
        and compare to the 505~ms / 6006~ms simulation values.

  \item \textbf{Fat-tree scale-out.} Extend the Mininet testbed from
        $k{=}4$ (16 hosts) to $k{=}6$ (54 hosts) and $k{=}8$
        (128 hosts). Measure orchestrator throughput (flows/s) and
        identify the saturation point.

  \item \textbf{Stronger adversarial evaluation.} Implement PGD and
        BIM attacks; add an adaptive adversary scenario (S9) where
        the attacker observes the feature thresholds and crafts
        traffic to stay below them. Measure model degradation.

  \item \textbf{Concept-drift monitor.} Implement a PSI
        (Population Stability Index) monitor on the 17 features;
        raise an alert when $\text{PSI} > 0.2$ and trigger a
        retraining request via DMaaP. This closes the adaptation
        loop missing from the current architecture.

  \item \textbf{Cross-dataset validation.} Apply the trained models
        (zero-shot) to BCCC-cPacket-Cloud-DDoS-2024. An accuracy
        above 85\% on that dataset would upgrade the thesis claim
        from ``in-distribution'' to ``cross-dataset generalising''.
\end{enumerate}

\medskip
\noindent\textbf{Closing remark.}
The central insight of this work is structural: reactive DDoS
mitigation has a hard latency floor equal to the VNF boot time, and
no orchestration tuning can lower that floor. Only \emph{prediction}
can --- by instantiating the VNF before the attack crosses the
detection threshold. \padonap{} operationalises this insight in an
ONAP-compliant, SLA-aware, and interpretable orchestration pipeline,
and demonstrates the 10.9$\times$ speed-up and 65~s lead-time
advantage on real evaluation scenarios.
